% ==========================================================
% 恰好两名女生相邻（捆绑与插空结合）
% 难度：★★ 中档
% ==========================================================

\newcommand{\qProbPermExactAdjStem}{%
三名男生和三名女生站成一排照相。男生甲与男生乙相邻，且三名女生中恰好有两名女生相邻。求不同站法数。
}

\newcommand{\qProbPermExactAdjAnswer}{%
$144$
}

\newcommand{\qProbPermExactAdjSolution}{%
\textbf{第一步：安排男生}\\
男生甲、乙必须相邻，把甲、乙捆绑成一个整体，\\
内部有 $A_2^2=2$ 种顺序。\\
把“甲乙整体”与另一名男生排列，有 $A_2^2=2$ 种。\\
因此男生部分共有 $2\times 2=4$ 种排列。\\
男生排好后形成三个空隙：\\
$\boxed{\text{空}}\quad \text{男生}\quad\boxed{\text{空}}\quad \text{男生}\quad\boxed{\text{空}}$。

\textbf{第二步：构造女生的“恰好相邻”}\\
从三名女生中选择两名并安排先后顺序，\\
组成相邻双人整体：$A_3^2=6$ 种。\\
此时女生分为两个单位：双人整体、独立女生。\\
为保证第三名女生不与双人整体连接，\\
这两个单位必须插入男生形成的两个不同空隙。\\
把两个不同单位插入三个不同空隙：$A_3^2=6$ 种。

\textbf{第三步：分步乘法}\\
总站法数为 $4\times 6\times 6 = 144$ 种。
}

\newcommand{\qProbPermExactAdjTeachingNotes}{%
\textbf{【避坑与边界说明】}
\begin{itemize}
    \item 只捆绑两名女生，不能保证“恰好两名女生相邻”，第三名女生可能贴在旁边。
    \item 女生双人整体与独立女生是两个不同单位，插入空隙时要考虑顺序。
    \item 两个女生单位必须进入不同空隙。
    \item 甲、乙捆绑后，内部仍有两种顺序。
\end{itemize}
}
